\chapter{Universality and complexity models}
A molecular machine can implement many programs without being a universal
computer. Conversely, a mathematically universal model can be physically
impractical. The first claim concerns which finite programs are available;
the second concerns an unbounded family of computations. Neither specifies
the preparation cost, failure probability or time of an experiment.

Chapter 17 exposed the storage hidden in reader positions. We now make
storage writable and reconstruct one tape step using two stacks. This gives
a concrete simulation invariant rather than a slogan about DNA's information
density. Our executable machine is a teaching model, not a molecular design.

\section{Start with an exact configuration}
A tape configuration is $(q,h,t)$: a control state $q$, an integer head
position $h$, and a map $t:\mathbb Z\to\{0,1,\_\}$ with finitely many
nonblank entries. An absent dictionary entry denotes blank. The tape is
bi-infinite; negative addresses are legal. A transition
\begin{equation}
 \delta(q,t(h))=(q',b,d),\qquad d\in\{-1,0,1\},
\end{equation}
writes $b$ at the old head, then moves by $d$ and enters $q'$. These
operations have an order. Moving before writing implements a different
machine unless the address used for the write is preserved explicitly.

\Plate{01-tape}{One carry transition reads 1, writes 0 at the same address,
then advances. The next read occurs at the new head. A tape diagram must
distinguish the control state, stored symbols and head position; an arrow
alone cannot specify the transition.}{fig:tape}

We distinguish three outcomes. \Term{Halt} means the designated halt state
was entered. \Term{Stuck} means a required transition was absent.
\Term{Exhausted} means the externally supplied execution budget ran out.
Exhaustion says nothing about whether more steps would eventually halt.
The interface starts in state \texttt{carry}, even for custom programs.
The companion validates every transition before running it.

\section{Build an incrementer, then test its contract}
Write binary input least-significant bit first. The word 110 denotes three,
not six. In state carry, a 1 becomes 0 and the head advances right; a 0
becomes 1 and halts; a blank becomes 1 and halts. The blank rule handles
overflow and also gives the empty word, interpreted as zero, an output of one.

\Listing{run}

For 111, four transitions produce 0001, representing eight:
\begin{center}\input{\Generated/trace.tex}\end{center}
Rows log the address \emph{before} the move. Tape span includes initial input
cells and all visited head positions, even blank ones. It is not the count of
nonblank entries, the dictionary's allocated bytes or the total history size.

Let $r$ be the number of leading 1s in this least-significant-first string.
The machine changes those $r$ bits to zero, then writes a 1 at position $r$.
The removed value is $\sum_{i=0}^{r-1}2^i=2^r-1$; the added value is
$2^r$. The net change is therefore exactly one. Higher positions remain
unchanged. This proves correctness for every finite input, including the
all-ones overflow case.

\Plate{03-growth}{Generated worst-case counts for all-ones inputs. Steps
and visited span both equal $n+1$ and therefore overlap. Equal plotted values
do not make the quantities interchangeable: one counts transitions, the
other counts tape addresses.}{fig:growth}

The exact step count is $r+1$, at most $n+1$. A word beginning with 0
halts after one transition even if the input is long. Reading the entire
input into a Python dictionary still costs work proportional to $n$.
An end-to-end runtime cannot discard that preparation just because the
abstract transition count is one.

\section{Replace the tape with two stacks}
Represent the cells strictly left of the head in stack $L$, with its top
nearest the head. Stack $R$ has the current cell at its top and the cells
to its right below it. Empty tails denote indefinitely many blanks. Python
list ends are stack tops. Initially $L$ is empty and $R$ is the reversed
input word.

After replacing the current top of $R$ with the written symbol, a right move
pops that symbol onto $L$. A left move pops the nearest left symbol onto
$R$, or pushes blank if $L$ is empty. A stationary move changes neither
stack boundary. The code below is the local operation; the companion's
\texttt{two\_stack} driver handles validation, budget and output decoding.
\Listing{stack_step}

\Plate{02-simulation}{A simulation invariant connects the tape and stack
representations. The right stack's top is the current cell. Encoding, taking
the stack step and decoding must agree with the tape transition, not merely
produce the same final answer on a few increment examples.}{fig:simulation}

\subsection{The invariant makes the simulation auditable}
At head $h$, the $i$th entry down from the top of $L$ represents
$t(h-1-i)$; the $i$th entry down from the top of $R$ represents $t(h+i)$.
For a right move, the written current symbol becomes the new nearest-left
symbol and the former right neighbor becomes current. For a left move,
the old nearest-left symbol becomes current and the written cell becomes
its immediate right neighbor. These are exactly the tape semantics.
They preserve the invariant by induction on the number of transitions.

Let $E$ encode a configuration, $D$ decode it and $F$ take a tape step.
With $G$ the corresponding stack step, the required equation is
\begin{equation}
 D(G(E(c)))=F(c).
\end{equation}
More generally a target may use $k(c)$ microsteps for one source step,
giving $D(G^{k(c)}(E(c)))=F(c)$. A simulation proof must specify the
intermediate states on which decoding is allowed and bound $k(c)$ when
making a complexity claim. Correct final answers alone do not establish this.

The tests compare the dictionary-tape and two-stack implementations on 500
seeded generated programs, including missing rules, left moves, blank writes
and bounded nonhalting loops. They compare status, trace, tape and span,
not only the incremented integer. This is finite differential evidence for
two independently represented machines; the invariant explains why the
representation works beyond those cases.

\section{What a universality claim must quantify}
\Term{Universality} requires a fixed interpreter that can simulate every
machine in a specified source class given an effective program encoding and
input. The incrementer's three rules are one program, not that interpreter.
An interpreter for encoded transition tables, with unbounded writable tape
or equivalent stacks and enough symbol encoding machinery, supplies the
general construction. A finite Python run on finite hardware realizes only
a bounded portion of this ideal model.

The two-stack construction demonstrates storage equivalence under its stated
operations. To carry that result into chemistry one would still need a
representation of arbitrarily extensible stacks, reliable top access,
transition control and a decoder. Giving a tube more distinct DNA species
does not by itself implement those operations. A molecular proposal should
state which part is a theorem about the model and which part is an
experiment on a bounded instance.

There is also a fundamental limit to testing. Suppose an algorithm decided
whether any encoded program halts on any input. Construct a program that
uses that decision on its own encoding, looping if the decision says halt
and halting if it says loop. Either answer contradicts the constructed
program's behavior. The contradiction concerns a general decider; it does
not prevent proving termination for this incrementer using a bounded carry
argument. A run budget is an engineering safeguard, not a halting oracle.

\section{Uniformity prevents answers from hiding in preparation}
For problem size $n$, let $D_n$ describe a device and let a finite constructor
produce $D_n$ from $n$. A meaningful cost claim accounts for the constructor,
device size, instance encoding, execution and readout. If the constructor
already solves the instance and places the answer into a reporter, the
subsequent one-step readout is not a one-step solution to the original problem.

\Plate{04-uniform}{A resource boundary includes device construction and
instance preparation, not just the reaction window. A red dashed path marks
answer-dependent preparation: it can conceal the very computation being
claimed. Uniformity must specify how the device family is generated.}{fig:uniform}

A useful ledger is
\begin{equation}
 C=(T_{\rm construct},S_{\rm device},T_{\rm encode},
 T_{\rm execute},V,N_{\rm species},T_{\rm readout}).
\end{equation}
These entries have different units and cannot be added without a model.
If preparation and execution overlap, wall-clock cost is a schedule, not
simply their sum. If a device is reused, amortization must state how many
instances share its construction cost. Even a reusable device cannot
amortize away per-instance encoding or detection.

Real-number models need an additional precision ledger. A single concentration
variable cannot serve as a free exact infinite-bit tape. The separation
between distinguishable levels, measurement noise and the required confidence
constrain usable information. A constant number of named chemical species
does not imply constant physical storage at arbitrary precision.

\section{Count opportunities for error}
Assume $T$ required events, each with error probability $\epsilon$, mutually
independent. Then the probability of no event error is
\begin{equation}
 P_0=(1-\epsilon)^T.
\end{equation}
Without independence, if every event's marginal error probability is at most
$\epsilon$, the union bound still gives $P_0\geq\max(0,1-T\epsilon)$.
Neither expression is automatically an output-correctness probability:
some event errors are harmless or repaired; omitted failure modes may spoil
the answer even when all counted events succeed.
\Listing{reliability}

\Plate{05-reliability}{Assigned error probability $\epsilon=0.001$ illustrates
the difference between an independent no-error model and a dependence-agnostic
union lower bound. These are calculated curves, not measured molecular
fidelities. A zero lower bound is uninformative, not a prediction of certain
failure.}{fig:reliability}

For $T=1000$, independence gives approximately 0.3677 whereas the union
lower bound is zero. For a target no-error probability at least $1-\delta$,
independence requires $\epsilon\leq1-(1-\delta)^{1/T}$.
For small $\delta$, this is roughly $\delta/T$. The union argument gives
the sufficient condition $\epsilon\leq\delta/T$ without independence.
Increasing problem size can therefore tighten fidelity requirements even
when the abstract computational model remains unchanged.

\section{A contemporary result and its boundary}
Stérin and colleagues' 2026 molecular-computer paper describes compilation of
finite-state machines into strand-displacement systems and a bounded
experimental demonstration \cite{sterin}. Its main text distinguishes kinetic
model conditions from experimentally enforced conditions. The relevant lesson
here is to separate compilation, kinetics and demonstrated instances. We do
not infer unbounded Turing universality, asymptotic advantage or arbitrary-length
experimental reliability from that result.

For a new architecture, ask for an explicit transition encoding, a simulation
invariant, construction and execution costs, a noise model, and a test that
could falsify the claimed mapping. A persuasive result need not solve every
layer at once; it must say which layer it establishes.

\section{Exercises with worked answers}
\begin{enumerate}
\item Increment 101. \textbf{Answer:} It represents five. The first 1
becomes 0, the next 0 becomes 1, producing 011, which represents six.
\item Increment the empty word. \textbf{Answer:} The blank rule writes 1
at address zero and halts in one step. The span is one cell.
\item Give a one-step long input. \textbf{Answer:} Any word beginning with
0; only its first bit changes. Loading the input is still a separate cost.
\item Run 111 with fuel three. \textbf{Answer:} The result is exhausted
after three writes; the fourth transition needed to write the overflow bit
has not happened. Exhausted must not be reported as halt.
\item Why retain visited blank addresses in span? \textbf{Answer:} Moving
through a blank still uses an addressable tape region. Sparse content and
visited extent are different resources.
\item Derive a left stack move. \textbf{Answer:} Write the current right
top, then push the left top onto the right stack. If the left stack is
empty, push blank. The new top is the new current cell.
\item Is constant stack-operation overhead constant physical time?
\textbf{Answer:} Only under the declared stack cost model. A chemical
implementation must justify access, transport and synchronization costs.
\item Why is the incrementer not universal? \textbf{Answer:} Its fixed
rules only perform increment. No arbitrary program encoding is interpreted.
\item Give a misleading constant-time claim. \textbf{Answer:} Solve an
instance during synthesis, encode the result into one reporter, and count
only reading that reporter. The preparation performs the omitted work.
\item Bound no-error probability for $T=100$, $\epsilon=0.001$.
\textbf{Answer:} The union lower bound is 0.9; independence gives about
0.9048. The latter needs an additional assumption.
\item Does a zero union bound predict failure? \textbf{Answer:} No. It
says this inequality no longer gives a positive guarantee.
\item Specify a universality evidence package. \textbf{Answer:} Define
the source class, program encoding, target states, step simulation, decoder
and unbounded resource assumptions; add separate construction costs and
finite experimental validation. Neither package substitutes for the other.
\end{enumerate}

\section{Vocabulary and the next question}
\Term{Simulation invariant}: a relation between representations preserved
by each simulated transition. \Term{Uniform construction}: an effective
procedure generating the device family under a stated cost bound.
\Term{Cost model}: the operations and resources assigned prices in an analysis.
\Term{Union bound}: an upper bound on any-error probability from a sum of
marginal error probabilities, without an independence assumption.

The next chapter begins the engineering question with toehold-mediated strand
displacement: how do binding, branch migration and release implement a
transition, and where can leakage break it? Universality alone cannot answer
that question.
